\documentclass[12pt]{article}
\usepackage[utf8]{inputenc}
\usepackage[T1]{fontenc}
\usepackage[margin=1in]{geometry}
\usepackage{setspace}
\usepackage{xcolor}
\usepackage[colorlinks=true, urlcolor=blue!60!black, linkcolor=blue!60!black]{hyperref}
\setlength{\parindent}{2em}
\setlength{\parskip}{1.25em}
\renewcommand{\baselinestretch}{1.0}
\begin{document}
\begin{center}
{\LARGE \textit{Sveen v. Melin}, 584 U.S. 811 (2018)\par}
\vspace{0.8em}
{\large Case Brief --- Four-Page Edition\par}
\end{center}
\vspace{1em}

\section{Memory Jogger}

Minnesota's revocation-on-divorce statute, applied retroactively, cut an ex-wife out as beneficiary of a life-insurance policy bought before the law passed. An 8--1 Court found no Contracts Clause violation because the law did not substantially impair the contract: it tracks likely intent, does what a divorce court could do, and can be undone with a simple form (\href{https://supreme.justia.com/cases/federal/us/584/16-1432/}{584 U.S.\ 811}).

\section{Detailed Case Facts}

In 1997 Mark Sveen bought a revocable life-insurance policy from MetLife and married Kaye Melin. In 1998 he named Melin primary beneficiary and his two adult children from a prior marriage, Ashley and Antone Sveen, contingent beneficiaries.

In 2002 Minnesota extended its revocation-on-divorce rule from wills to ``will substitutes'' such as life insurance. The statute provides that divorce ``revokes any revocable \ldots{} beneficiary designation \ldots{} made by an individual to the individual's former spouse,'' treating the ex-spouse as having predeceased the policyholder unless the policyholder re-designates her or the decree provides otherwise. Mark and Melin divorced in 2007; the decree did not mention the policy, and Mark never changed the designation. When he died in 2011, both Melin and his children claimed the proceeds.

\section{Procedural History}

MetLife filed an interpleader action in the District of Minnesota. The district court granted summary judgment to the children, holding that retroactive application of the statute did not violate the Contracts Clause. The Eighth Circuit reversed, holding that retroactive application disrupted the policyholder's expectations. Given a split among the Eighth, Ninth, and Tenth Circuits and the Pennsylvania and Colorado Supreme Courts, the Court granted certiorari, heard argument on March 19, 2018, and reversed on June 11, 2018.

\section{Judicial Votes}

Reversed and remanded, \textbf{8--1}.
\begin{itemize}
\item \textbf{Majority:} Justice Kagan, joined by Chief Justice Roberts and Justices Kennedy, Thomas, Ginsburg, Breyer, Alito, and Sotomayor.
\item \textbf{Dissenting:} Justice Gorsuch.
\end{itemize}

\section{Holding}

Retroactive application of Minnesota's revocation-on-divorce statute does not violate the Contracts Clause of Article I, Section 10, because the law does not cross the step-one threshold of ``substantial impairment'' of a contractual relationship. The Court therefore did not reach step two.

\section{Analysis of Opinions}

\textbf{Majority (Kagan).} Not every state law affecting existing contracts violates the Clause (\textit{El Paso v.\ Simmons}). Under the two-step framework, a court first asks whether the law ``operated as a substantial impairment of a contractual relationship'' (\textit{Allied Structural Steel}), considering how far it undermines the bargain, interferes with reasonable expectations, and prevents the party from safeguarding his rights; only then does it ask whether the law is a reasonable and appropriate means to ``a significant and legitimate public purpose'' (\textit{Energy Reserves}). Declining in footnote 2 to revisit the test, Kagan found three features decisive. \emph{Intent:} the statute rests on the presumption that ``the average Joe does not want his ex inheriting what he leaves behind,'' so it supports the contractual scheme. \emph{Expectations:} divorce courts have long had power to reallocate life insurance, so no policyholder could expect a spousal designation to survive divorce untouched. \emph{Reversibility:} the default can be undone ``with the stroke of a pen.'' Two centuries of precedent, from \textit{Jackson v.\ Lamphire} (1830) to \textit{Texaco, Inc.\ v.\ Short} (1982), upheld retroactive recording and filing laws that extinguished whole property rights, while this law merely redirects proceeds. Following \textit{Texaco}, the Court rejected a distinction between altering rights and regulating remedies.

\textbf{Gorsuch (dissenting).} Relying on the text and \textit{Sturges v.\ Crowninshield}, Gorsuch argued that the Clause categorically forbids any retroactive impairment and that critics of the post-\textit{Blaisdell} balancing test ``deserve a thoughtful reply.'' Even under the test, naming the beneficiary is the ``whole point'' of the contract; legislatures cannot guess every divorcing spouse's intent, since some keep ex-spouses as beneficiaries for child support; a blanket rewrite differs from a decree dividing specific property; and a default that catches inattentive policyholders still destroys their contract. Less intrusive options, such as having divorce courts notify spouses to review their forms, were available.

\section{Examples: Future Applications}

\begin{enumerate}
\item \textbf{Same-Side: Bank Accounts.} Extending revocation-on-divorce to existing payable-on-death bank accounts satisfies all three \textit{Sveen} factors.
\item \textbf{Same-Side: Mineral Interests.} A law lapsing dormant mineral interests unless owners file a short statement within two years is the kind upheld in \textit{Texaco v.\ Short}.
\item \textbf{Opposite-Side: Voiding Non-Competes.} Retroactively voiding all existing non-competes and lease guaranties with no opt-out fails every step-one factor, so step two applies.
\item \textbf{Opposite-Side: Repealing a Bond Covenant.} A State's repeal of its own toll-revenue covenant to bondholders destroys a core term and faces non-deferential step-two review under \textit{United States Trust Co.\ v.\ New Jersey} (1977).
\item \textbf{Fence-Sitter: Hard-to-Undo Revocation.} If re-designation requires a court petition, fee, and hearing, the third factor flips, and \textit{Sveen} did not say whether the first two suffice alone.
\end{enumerate}

\section{Critique}

\textbf{Notice and the unwary.} The ``stroke of a pen'' rationale assumes policyholders know the legislature changed their existing contracts. An inattentive policyholder who relied on the law as it stood has no reason to file a new form to keep an existing designation.

\textbf{Collapse of step one.} By folding probable intent and legislative reasonableness into the substantial-impairment threshold, \textit{Sveen} lets courts dismiss Contracts Clause claims without ever examining tailoring at step two. Note too that the Clause binds only the States; retroactive federal laws are reviewed under Fifth Amendment due process or the Takings Clause.

\section{Key Quotations}

\begin{quote}
``In answering that question, the Court has considered the extent to which the law undermines the contractual bargain, interferes with a party's reasonable expectations, and prevents the party from safeguarding or reinstating his rights.''
\end{quote}
\noindent --- Justice Kagan, for the Court, 584 U.S.\ at 819.

\begin{quote}
``First, the statute is designed to reflect a policyholder's intent---and so to support, rather than impair, the contractual scheme. Second, the law is unlikely to disturb any policyholder's expectations because it does no more than a divorce court could always have done. And third, the statute supplies a mere default rule, which the policyholder can undo in a moment.''
\end{quote}
\noindent --- Justice Kagan, for the Court, 584 U.S.\ at 820.

\begin{quote}
``In cases going back to the 1800s, this Court has held that laws imposing such minimal paperwork burdens do not violate the Contracts Clause.''
\end{quote}
\noindent --- Justice Kagan, for the Court, 584 U.S.\ at 822.

\end{document}
